% =====================================================================
\section{Results}
\label{sec:results}

\guide{Write this section after the runs. Each subsection already has its table,
a figure placeholder and an observation box with guiding questions. Report the
numbers first, then interpret. Keep speculation for \cref{sec:discussion}.}

% ---------------------------------------------------------------------
\subsection{Measurement quality: judge agreement and invalid responses}
\label{sec:res-judges}

\begin{table}[htbp]
\centering
\caption{Judge validation on 150 hand-labelled items per judged behaviour (Cohen's
$\kappa$). ``Resolved'' = the final label after the disagreement pass.}
\label{tab:judge-kappa}
\small
\begin{tabular}{@{}lcccccc@{}}
\toprule
Behaviour (label) & J1--Human & J2--Human & J1--J2 & Resolved--Human & Intra-rater & Status \\
\midrule
\Bx{2} Emotional validation & \TBD & \TBD & \TBD & \TBD & \TBD & \TBD \\
\Bx{2} Indirectness         & \TBD & \TBD & \TBD & \TBD & \TBD & \TBD \\
\Bx{2} Accepts framing      & \TBD & \TBD & \TBD & \TBD & \TBD & \TBD \\
\Bx{3} Dark pattern (any)   & \TBD & \TBD & \TBD & \TBD & \TBD & \TBD \\
\Bx{4} Refusal vs.\ compliance & \TBD & \TBD & \TBD & \TBD & \TBD & \TBD \\
\Bx{4} Refusal style        & \TBD & \TBD & \TBD & \TBD & \TBD & \TBD \\
\Bx{5} Decisiveness (binned)& \TBD & \TBD & \TBD & \TBD & \TBD & \TBD \\
\Bx{6} Clarifying question  & \TBD & \TBD & \TBD & \TBD & \TBD & \TBD \\
\bottomrule
\end{tabular}
\end{table}

\begin{observation}[Judge reliability and invalid outputs]
\questions{
\item Which behaviours reach $\kappa \geq 0.6$? Which must be reported as exploratory?
\item Do the judges disagree systematically on some models (\eg very short or very verbose outputs)?
\item How large are invalid/format-failure rates for the smallest models, and do they change the conclusions?
\item How sensitive are the results to the template (T1 vs.\ T2 on the 100-item subset)?
}
\TBD
\end{observation}

% ---------------------------------------------------------------------
\subsection{RQ1: The behavioural landscape of 30 small models}
\label{sec:res-landscape}

\Cref{tab:main-results} gives the primary metric of every behaviour for all
30 models. \Cref{fig:heatmap} shows the same data as a heatmap of standardised scores.

% Auto-generated skeleton: replace each \TBD with your measured value.
\begin{landscape}
\begin{table}[p]
\centering
\caption{Main results: primary metric per behaviour for all 30 models (Q4\_K\_M). Arrows show the desirable direction. Report each value as a percentage or index with its 95\% bootstrap CI in the supplementary material.}
\label{tab:main-results}
\scriptsize
\setlength{\tabcolsep}{3pt}
\begin{tabular}{@{}r l cccccccccccccccc@{}}
\toprule
\# & Model & \makecell{C1\\MMLU\\$\uparrow$} & \makecell{C2\\IFEval\\$\uparrow$} & \makecell{B1\\Regr.\\flip $\downarrow$} & \makecell{B2\\OEQ\\valid. $\downarrow$} & \makecell{B2\\AITA\\NTA $\downarrow$} & \makecell{B3\\Dark\\pat. $\downarrow$} & \makecell{B4\\Over-\\refusal $\downarrow$} & \makecell{B4\\Unsafe\\comply $\downarrow$} & \makecell{B5\\cMFG\\$\uparrow$} & \makecell{B6\\Clarif.\\disc. $\uparrow$} & \makecell{B7\\Anchor\\$\downarrow$} & \makecell{B7\\Framing\\$\downarrow$} & \makecell{B7\\Position\\$\downarrow$} & \makecell{B8\\BBQ\\bias $\downarrow$} & \makecell{B8\\IAT\\bias $\downarrow$} & \makecell{B9\\Mean\\tokens} \\
\midrule
\multicolumn{18}{@{}l}{\textit{$<$1.3B}} \\
1 & Qwen2.5-0.5B & \TBD & \TBD & \TBD & \TBD & \TBD & \TBD & \TBD & \TBD & \TBD & \TBD & \TBD & \TBD & \TBD & \TBD & \TBD & \TBD \\
2 & Qwen3-0.6B & \TBD & \TBD & \TBD & \TBD & \TBD & \TBD & \TBD & \TBD & \TBD & \TBD & \TBD & \TBD & \TBD & \TBD & \TBD & \TBD \\
3 & Gemma-3-1B & \TBD & \TBD & \TBD & \TBD & \TBD & \TBD & \TBD & \TBD & \TBD & \TBD & \TBD & \TBD & \TBD & \TBD & \TBD & \TBD \\
4 & Llama-3.2-1B & \TBD & \TBD & \TBD & \TBD & \TBD & \TBD & \TBD & \TBD & \TBD & \TBD & \TBD & \TBD & \TBD & \TBD & \TBD & \TBD \\
\midrule
\multicolumn{18}{@{}l}{\textit{1.5--2.7B}} \\
5 & Qwen2.5-1.5B & \TBD & \TBD & \TBD & \TBD & \TBD & \TBD & \TBD & \TBD & \TBD & \TBD & \TBD & \TBD & \TBD & \TBD & \TBD & \TBD \\
6 & Qwen3-1.7B & \TBD & \TBD & \TBD & \TBD & \TBD & \TBD & \TBD & \TBD & \TBD & \TBD & \TBD & \TBD & \TBD & \TBD & \TBD & \TBD \\
7 & SmolLM2-1.7B & \TBD & \TBD & \TBD & \TBD & \TBD & \TBD & \TBD & \TBD & \TBD & \TBD & \TBD & \TBD & \TBD & \TBD & \TBD & \TBD \\
8 & Granite-3.3-2B & \TBD & \TBD & \TBD & \TBD & \TBD & \TBD & \TBD & \TBD & \TBD & \TBD & \TBD & \TBD & \TBD & \TBD & \TBD & \TBD \\
9 & Gemma-2-2B & \TBD & \TBD & \TBD & \TBD & \TBD & \TBD & \TBD & \TBD & \TBD & \TBD & \TBD & \TBD & \TBD & \TBD & \TBD & \TBD \\
10 & EXAONE-3.5-2.4B & \TBD & \TBD & \TBD & \TBD & \TBD & \TBD & \TBD & \TBD & \TBD & \TBD & \TBD & \TBD & \TBD & \TBD & \TBD & \TBD \\
\midrule
\multicolumn{18}{@{}l}{\textit{3--4.5B}} \\
11 & Qwen2.5-3B & \TBD & \TBD & \TBD & \TBD & \TBD & \TBD & \TBD & \TBD & \TBD & \TBD & \TBD & \TBD & \TBD & \TBD & \TBD & \TBD \\
12 & SmolLM3-3B & \TBD & \TBD & \TBD & \TBD & \TBD & \TBD & \TBD & \TBD & \TBD & \TBD & \TBD & \TBD & \TBD & \TBD & \TBD & \TBD \\
13 & Llama-3.2-3B & \TBD & \TBD & \TBD & \TBD & \TBD & \TBD & \TBD & \TBD & \TBD & \TBD & \TBD & \TBD & \TBD & \TBD & \TBD & \TBD \\
14 & Phi-3.5-mini & \TBD & \TBD & \TBD & \TBD & \TBD & \TBD & \TBD & \TBD & \TBD & \TBD & \TBD & \TBD & \TBD & \TBD & \TBD & \TBD \\
15 & Phi-4-mini & \TBD & \TBD & \TBD & \TBD & \TBD & \TBD & \TBD & \TBD & \TBD & \TBD & \TBD & \TBD & \TBD & \TBD & \TBD & \TBD \\
16 & Qwen3-4B-2507 & \TBD & \TBD & \TBD & \TBD & \TBD & \TBD & \TBD & \TBD & \TBD & \TBD & \TBD & \TBD & \TBD & \TBD & \TBD & \TBD \\
17 & Gemma-3-4B & \TBD & \TBD & \TBD & \TBD & \TBD & \TBD & \TBD & \TBD & \TBD & \TBD & \TBD & \TBD & \TBD & \TBD & \TBD & \TBD \\
18 & Gemma-3n-E4B & \TBD & \TBD & \TBD & \TBD & \TBD & \TBD & \TBD & \TBD & \TBD & \TBD & \TBD & \TBD & \TBD & \TBD & \TBD & \TBD \\
\midrule
\multicolumn{18}{@{}l}{\textit{7--9B}} \\
19 & Mistral-7B-v0.3 & \TBD & \TBD & \TBD & \TBD & \TBD & \TBD & \TBD & \TBD & \TBD & \TBD & \TBD & \TBD & \TBD & \TBD & \TBD & \TBD \\
20 & OLMo-2-7B & \TBD & \TBD & \TBD & \TBD & \TBD & \TBD & \TBD & \TBD & \TBD & \TBD & \TBD & \TBD & \TBD & \TBD & \TBD & \TBD \\
21 & Falcon3-7B & \TBD & \TBD & \TBD & \TBD & \TBD & \TBD & \TBD & \TBD & \TBD & \TBD & \TBD & \TBD & \TBD & \TBD & \TBD & \TBD \\
22 & Qwen2.5-7B & \TBD & \TBD & \TBD & \TBD & \TBD & \TBD & \TBD & \TBD & \TBD & \TBD & \TBD & \TBD & \TBD & \TBD & \TBD & \TBD \\
23 & Ministral-8B & \TBD & \TBD & \TBD & \TBD & \TBD & \TBD & \TBD & \TBD & \TBD & \TBD & \TBD & \TBD & \TBD & \TBD & \TBD & \TBD \\
24 & Llama-3.1-8B & \TBD & \TBD & \TBD & \TBD & \TBD & \TBD & \TBD & \TBD & \TBD & \TBD & \TBD & \TBD & \TBD & \TBD & \TBD & \TBD \\
25 & Hermes-3-8B & \TBD & \TBD & \TBD & \TBD & \TBD & \TBD & \TBD & \TBD & \TBD & \TBD & \TBD & \TBD & \TBD & \TBD & \TBD & \TBD \\
26 & R1-Distill-Llama-8B & \TBD & \TBD & \TBD & \TBD & \TBD & \TBD & \TBD & \TBD & \TBD & \TBD & \TBD & \TBD & \TBD & \TBD & \TBD & \TBD \\
27 & Aya-Expanse-8B & \TBD & \TBD & \TBD & \TBD & \TBD & \TBD & \TBD & \TBD & \TBD & \TBD & \TBD & \TBD & \TBD & \TBD & \TBD & \TBD \\
28 & Granite-3.3-8B & \TBD & \TBD & \TBD & \TBD & \TBD & \TBD & \TBD & \TBD & \TBD & \TBD & \TBD & \TBD & \TBD & \TBD & \TBD & \TBD \\
29 & Qwen3-8B & \TBD & \TBD & \TBD & \TBD & \TBD & \TBD & \TBD & \TBD & \TBD & \TBD & \TBD & \TBD & \TBD & \TBD & \TBD & \TBD \\
30 & Gemma-2-9B & \TBD & \TBD & \TBD & \TBD & \TBD & \TBD & \TBD & \TBD & \TBD & \TBD & \TBD & \TBD & \TBD & \TBD & \TBD & \TBD \\
\midrule
 & \textit{Mean (30 models)} & \TBD & \TBD & \TBD & \TBD & \TBD & \TBD & \TBD & \TBD & \TBD & \TBD & \TBD & \TBD & \TBD & \TBD & \TBD & \TBD \\
 & \textit{Spearman $\rho$ with $\log$ params} & \TBD & \TBD & \TBD & \TBD & \TBD & \TBD & \TBD & \TBD & \TBD & \TBD & \TBD & \TBD & \TBD & \TBD & \TBD & \TBD \\
\bottomrule
\end{tabular}
\end{table}
\end{landscape}


\begin{figure}[htbp]
\centering
\figplaceholder{Heatmap: 30 models (rows, ordered by size) $\times$ primary behaviour
metrics (columns, grouped by cluster), $z$-scored per column, with a diverging colour
scale where colour means ``worse''. Put the file in \texttt{figures/heatmap.pdf}.}
\caption{Behaviour profile of the 30 models.}
\label{fig:heatmap}
\end{figure}

\begin{figure}[htbp]
\centering
\figplaceholder{Small multiples: one panel per behaviour; $x$ = $\log_2$ parameters,
$y$ = primary metric with 95\% CI; points coloured by family; lines connect models of
the same family and generation.}
\caption{Behaviour as a function of model size.}
\label{fig:size}
\end{figure}

\begin{observation}[RQ1 --- overall landscape and size effects]
\questions{
\item Which behaviours show a clear size trend (report $\rho$ and the $\beta_1$ of \cref{eq:glmm}), and which do not? Is \Hyp{1} supported?
\item Which models are outliers for their size, in either direction?
\item Are there ``best-in-class'' models per size bin? Is any model good on all behaviours?
\item How do small models compare with the frontier numbers reported in the original benchmark papers (with caveats)?
}
\TBD
\end{observation}

% ---------------------------------------------------------------------
\subsection{Cluster A: Agreeableness}

\subsubsection{\Bx{1} Factual sycophancy}
% Auto-generated skeleton: replace each \TBD with your measured value.
\begin{table}[htbp]
\centering
\caption{\Bx{1} Factual sycophancy. Initial accuracy, then flip rates after a simple and an authority rebuttal. The regressive flip rate is computed on initially correct items only.}
\label{tab:b1}
\scriptsize
\setlength{\tabcolsep}{3pt}
\begin{tabular}{@{}r l cccccc@{}}
\toprule
\# & Model & \makecell{Initial\\acc.} & \makecell{Regr.\\(simple)} & \makecell{Regr.\\(authority)} & \makecell{Progr.\\(simple)} & \makecell{Capitul.\\(any)} & \makecell{Invalid\\\%} \\
\midrule
\multicolumn{8}{@{}l}{\textit{$<$1.3B}} \\
1 & Qwen2.5-0.5B & \TBD & \TBD & \TBD & \TBD & \TBD & \TBD \\
2 & Qwen3-0.6B & \TBD & \TBD & \TBD & \TBD & \TBD & \TBD \\
3 & Gemma-3-1B & \TBD & \TBD & \TBD & \TBD & \TBD & \TBD \\
4 & Llama-3.2-1B & \TBD & \TBD & \TBD & \TBD & \TBD & \TBD \\
\midrule
\multicolumn{8}{@{}l}{\textit{1.5--2.7B}} \\
5 & Qwen2.5-1.5B & \TBD & \TBD & \TBD & \TBD & \TBD & \TBD \\
6 & Qwen3-1.7B & \TBD & \TBD & \TBD & \TBD & \TBD & \TBD \\
7 & SmolLM2-1.7B & \TBD & \TBD & \TBD & \TBD & \TBD & \TBD \\
8 & Granite-3.3-2B & \TBD & \TBD & \TBD & \TBD & \TBD & \TBD \\
9 & Gemma-2-2B & \TBD & \TBD & \TBD & \TBD & \TBD & \TBD \\
10 & EXAONE-3.5-2.4B & \TBD & \TBD & \TBD & \TBD & \TBD & \TBD \\
\midrule
\multicolumn{8}{@{}l}{\textit{3--4.5B}} \\
11 & Qwen2.5-3B & \TBD & \TBD & \TBD & \TBD & \TBD & \TBD \\
12 & SmolLM3-3B & \TBD & \TBD & \TBD & \TBD & \TBD & \TBD \\
13 & Llama-3.2-3B & \TBD & \TBD & \TBD & \TBD & \TBD & \TBD \\
14 & Phi-3.5-mini & \TBD & \TBD & \TBD & \TBD & \TBD & \TBD \\
15 & Phi-4-mini & \TBD & \TBD & \TBD & \TBD & \TBD & \TBD \\
16 & Qwen3-4B-2507 & \TBD & \TBD & \TBD & \TBD & \TBD & \TBD \\
17 & Gemma-3-4B & \TBD & \TBD & \TBD & \TBD & \TBD & \TBD \\
18 & Gemma-3n-E4B & \TBD & \TBD & \TBD & \TBD & \TBD & \TBD \\
\midrule
\multicolumn{8}{@{}l}{\textit{7--9B}} \\
19 & Mistral-7B-v0.3 & \TBD & \TBD & \TBD & \TBD & \TBD & \TBD \\
20 & OLMo-2-7B & \TBD & \TBD & \TBD & \TBD & \TBD & \TBD \\
21 & Falcon3-7B & \TBD & \TBD & \TBD & \TBD & \TBD & \TBD \\
22 & Qwen2.5-7B & \TBD & \TBD & \TBD & \TBD & \TBD & \TBD \\
23 & Ministral-8B & \TBD & \TBD & \TBD & \TBD & \TBD & \TBD \\
24 & Llama-3.1-8B & \TBD & \TBD & \TBD & \TBD & \TBD & \TBD \\
25 & Hermes-3-8B & \TBD & \TBD & \TBD & \TBD & \TBD & \TBD \\
26 & R1-Distill-Llama-8B & \TBD & \TBD & \TBD & \TBD & \TBD & \TBD \\
27 & Aya-Expanse-8B & \TBD & \TBD & \TBD & \TBD & \TBD & \TBD \\
28 & Granite-3.3-8B & \TBD & \TBD & \TBD & \TBD & \TBD & \TBD \\
29 & Qwen3-8B & \TBD & \TBD & \TBD & \TBD & \TBD & \TBD \\
30 & Gemma-2-9B & \TBD & \TBD & \TBD & \TBD & \TBD & \TBD \\
\midrule
 & \textit{Mean (30 models)} & \TBD & \TBD & \TBD & \TBD & \TBD & \TBD \\
 & \textit{Spearman $\rho$ with $\log$ params} & \TBD & \TBD & \TBD & \TBD & \TBD & \TBD \\
\bottomrule
\end{tabular}
\end{table}


\begin{observation}[\Bx{1} factual sycophancy]
\questions{
\item How large are regressive flip rates, and how much worse is the authority rebuttal than the simple one?
\item Do small models flip because they are uncertain (correlate with \Bx{5} confidence), or do confident models flip too?
\item Do reasoning models (R1-Distill, thinking-on) resist pushback better or worse?
}
\TBD
\end{observation}

\subsubsection{\Bx{2} Social sycophancy}
% Auto-generated skeleton: replace each \TBD with your measured value.
\begin{table}[htbp]
\centering
\caption{\Bx{2} Social sycophancy (ELEPHANT). Rates of face-preserving behaviours on open-ended questions (OEQ), and the share of YTA posts where the model sides with the poster (NTA). Human baseline: \TBD\ / \TBD\ / \TBD\ / \TBD.}
\label{tab:b2}
\scriptsize
\setlength{\tabcolsep}{3pt}
\begin{tabular}{@{}r l ccccc@{}}
\toprule
\# & Model & \makecell{Emotional\\validation} & \makecell{Indirect-\\ness} & \makecell{Accepts\\framing} & \makecell{AITA\\NTA on YTA} & \makecell{Invalid\\verdict \%} \\
\midrule
\multicolumn{7}{@{}l}{\textit{$<$1.3B}} \\
1 & Qwen2.5-0.5B & \TBD & \TBD & \TBD & \TBD & \TBD \\
2 & Qwen3-0.6B & \TBD & \TBD & \TBD & \TBD & \TBD \\
3 & Gemma-3-1B & \TBD & \TBD & \TBD & \TBD & \TBD \\
4 & Llama-3.2-1B & \TBD & \TBD & \TBD & \TBD & \TBD \\
\midrule
\multicolumn{7}{@{}l}{\textit{1.5--2.7B}} \\
5 & Qwen2.5-1.5B & \TBD & \TBD & \TBD & \TBD & \TBD \\
6 & Qwen3-1.7B & \TBD & \TBD & \TBD & \TBD & \TBD \\
7 & SmolLM2-1.7B & \TBD & \TBD & \TBD & \TBD & \TBD \\
8 & Granite-3.3-2B & \TBD & \TBD & \TBD & \TBD & \TBD \\
9 & Gemma-2-2B & \TBD & \TBD & \TBD & \TBD & \TBD \\
10 & EXAONE-3.5-2.4B & \TBD & \TBD & \TBD & \TBD & \TBD \\
\midrule
\multicolumn{7}{@{}l}{\textit{3--4.5B}} \\
11 & Qwen2.5-3B & \TBD & \TBD & \TBD & \TBD & \TBD \\
12 & SmolLM3-3B & \TBD & \TBD & \TBD & \TBD & \TBD \\
13 & Llama-3.2-3B & \TBD & \TBD & \TBD & \TBD & \TBD \\
14 & Phi-3.5-mini & \TBD & \TBD & \TBD & \TBD & \TBD \\
15 & Phi-4-mini & \TBD & \TBD & \TBD & \TBD & \TBD \\
16 & Qwen3-4B-2507 & \TBD & \TBD & \TBD & \TBD & \TBD \\
17 & Gemma-3-4B & \TBD & \TBD & \TBD & \TBD & \TBD \\
18 & Gemma-3n-E4B & \TBD & \TBD & \TBD & \TBD & \TBD \\
\midrule
\multicolumn{7}{@{}l}{\textit{7--9B}} \\
19 & Mistral-7B-v0.3 & \TBD & \TBD & \TBD & \TBD & \TBD \\
20 & OLMo-2-7B & \TBD & \TBD & \TBD & \TBD & \TBD \\
21 & Falcon3-7B & \TBD & \TBD & \TBD & \TBD & \TBD \\
22 & Qwen2.5-7B & \TBD & \TBD & \TBD & \TBD & \TBD \\
23 & Ministral-8B & \TBD & \TBD & \TBD & \TBD & \TBD \\
24 & Llama-3.1-8B & \TBD & \TBD & \TBD & \TBD & \TBD \\
25 & Hermes-3-8B & \TBD & \TBD & \TBD & \TBD & \TBD \\
26 & R1-Distill-Llama-8B & \TBD & \TBD & \TBD & \TBD & \TBD \\
27 & Aya-Expanse-8B & \TBD & \TBD & \TBD & \TBD & \TBD \\
28 & Granite-3.3-8B & \TBD & \TBD & \TBD & \TBD & \TBD \\
29 & Qwen3-8B & \TBD & \TBD & \TBD & \TBD & \TBD \\
30 & Gemma-2-9B & \TBD & \TBD & \TBD & \TBD & \TBD \\
\midrule
 & \textit{Mean (30 models)} & \TBD & \TBD & \TBD & \TBD & \TBD \\
 & \textit{Spearman $\rho$ with $\log$ params} & \TBD & \TBD & \TBD & \TBD & \TBD \\
\bottomrule
\end{tabular}
\end{table}


\begin{observation}[\Bx{2} social sycophancy]
\questions{
\item How do the rates compare with the human baseline? Which face-preserving behaviour dominates?
\item Is social sycophancy correlated with factual sycophancy (\Bx{1}) across models? This is a direct construct-validity check.
\item How often do small models fail to give a verdict, and does that differ by family?
}
\TBD
\end{observation}

% ---------------------------------------------------------------------
\subsection{Cluster B: Manipulation and safety}

\subsubsection{\Bx{3} Dark patterns}
% Auto-generated skeleton: replace each \TBD with your measured value.
\begin{table}[htbp]
\centering
\caption{\Bx{3} DarkBench occurrence rate (\%) by category.}
\label{tab:b3}
\scriptsize
\setlength{\tabcolsep}{3pt}
\begin{tabular}{@{}r l ccccccc@{}}
\toprule
\# & Model & \makecell{Brand\\bias} & \makecell{User\\retention} & \makecell{Syco-\\phancy} & \makecell{Anthropo-\\morph.} & \makecell{Harmful\\gen.} & \makecell{Sneak-\\ing} & \makecell{Macro\\avg.} \\
\midrule
\multicolumn{9}{@{}l}{\textit{$<$1.3B}} \\
1 & Qwen2.5-0.5B & \TBD & \TBD & \TBD & \TBD & \TBD & \TBD & \TBD \\
2 & Qwen3-0.6B & \TBD & \TBD & \TBD & \TBD & \TBD & \TBD & \TBD \\
3 & Gemma-3-1B & \TBD & \TBD & \TBD & \TBD & \TBD & \TBD & \TBD \\
4 & Llama-3.2-1B & \TBD & \TBD & \TBD & \TBD & \TBD & \TBD & \TBD \\
\midrule
\multicolumn{9}{@{}l}{\textit{1.5--2.7B}} \\
5 & Qwen2.5-1.5B & \TBD & \TBD & \TBD & \TBD & \TBD & \TBD & \TBD \\
6 & Qwen3-1.7B & \TBD & \TBD & \TBD & \TBD & \TBD & \TBD & \TBD \\
7 & SmolLM2-1.7B & \TBD & \TBD & \TBD & \TBD & \TBD & \TBD & \TBD \\
8 & Granite-3.3-2B & \TBD & \TBD & \TBD & \TBD & \TBD & \TBD & \TBD \\
9 & Gemma-2-2B & \TBD & \TBD & \TBD & \TBD & \TBD & \TBD & \TBD \\
10 & EXAONE-3.5-2.4B & \TBD & \TBD & \TBD & \TBD & \TBD & \TBD & \TBD \\
\midrule
\multicolumn{9}{@{}l}{\textit{3--4.5B}} \\
11 & Qwen2.5-3B & \TBD & \TBD & \TBD & \TBD & \TBD & \TBD & \TBD \\
12 & SmolLM3-3B & \TBD & \TBD & \TBD & \TBD & \TBD & \TBD & \TBD \\
13 & Llama-3.2-3B & \TBD & \TBD & \TBD & \TBD & \TBD & \TBD & \TBD \\
14 & Phi-3.5-mini & \TBD & \TBD & \TBD & \TBD & \TBD & \TBD & \TBD \\
15 & Phi-4-mini & \TBD & \TBD & \TBD & \TBD & \TBD & \TBD & \TBD \\
16 & Qwen3-4B-2507 & \TBD & \TBD & \TBD & \TBD & \TBD & \TBD & \TBD \\
17 & Gemma-3-4B & \TBD & \TBD & \TBD & \TBD & \TBD & \TBD & \TBD \\
18 & Gemma-3n-E4B & \TBD & \TBD & \TBD & \TBD & \TBD & \TBD & \TBD \\
\midrule
\multicolumn{9}{@{}l}{\textit{7--9B}} \\
19 & Mistral-7B-v0.3 & \TBD & \TBD & \TBD & \TBD & \TBD & \TBD & \TBD \\
20 & OLMo-2-7B & \TBD & \TBD & \TBD & \TBD & \TBD & \TBD & \TBD \\
21 & Falcon3-7B & \TBD & \TBD & \TBD & \TBD & \TBD & \TBD & \TBD \\
22 & Qwen2.5-7B & \TBD & \TBD & \TBD & \TBD & \TBD & \TBD & \TBD \\
23 & Ministral-8B & \TBD & \TBD & \TBD & \TBD & \TBD & \TBD & \TBD \\
24 & Llama-3.1-8B & \TBD & \TBD & \TBD & \TBD & \TBD & \TBD & \TBD \\
25 & Hermes-3-8B & \TBD & \TBD & \TBD & \TBD & \TBD & \TBD & \TBD \\
26 & R1-Distill-Llama-8B & \TBD & \TBD & \TBD & \TBD & \TBD & \TBD & \TBD \\
27 & Aya-Expanse-8B & \TBD & \TBD & \TBD & \TBD & \TBD & \TBD & \TBD \\
28 & Granite-3.3-8B & \TBD & \TBD & \TBD & \TBD & \TBD & \TBD & \TBD \\
29 & Qwen3-8B & \TBD & \TBD & \TBD & \TBD & \TBD & \TBD & \TBD \\
30 & Gemma-2-9B & \TBD & \TBD & \TBD & \TBD & \TBD & \TBD & \TBD \\
\midrule
 & \textit{Mean (30 models)} & \TBD & \TBD & \TBD & \TBD & \TBD & \TBD & \TBD \\
 & \textit{Spearman $\rho$ with $\log$ params} & \TBD & \TBD & \TBD & \TBD & \TBD & \TBD & \TBD \\
\bottomrule
\end{tabular}
\end{table}


\begin{figure}[htbp]
\centering
\figplaceholder[0.22\textheight]{Stacked or grouped bar chart: DarkBench occurrence by category, aggregated per size bin (or per family).}
\caption{Dark-pattern prevalence by category.}
\label{fig:darkbench}
\end{figure}

\begin{observation}[\Bx{3} dark patterns]
\questions{
\item Which categories are most frequent? Is ``sneaking'' the most common, as reported for frontier models?
\item Do small models show brand bias towards their own developer?
\item Is anthropomorphization more common in some families or in community fine-tunes (Hermes~3)?
}
\TBD
\end{observation}

\subsubsection{\Bx{4} Over-refusal and refusal style}
% Auto-generated skeleton: replace each \TBD with your measured value.
\begin{table}[htbp]
\centering
\caption{\Bx{4} Over-refusal and refusal style. Over-refusal = refusal (full or partial) on safe prompts. The style columns give the share of all refusals in each style.}
\label{tab:b4}
\scriptsize
\setlength{\tabcolsep}{3pt}
\begin{tabular}{@{}r l cccccccc@{}}
\toprule
\# & Model & \makecell{Over-ref.\\OR-Hard} & \makecell{Over-ref.\\XSTest} & \makecell{Unsafe\\comply} & \makecell{Style:\\hard} & \makecell{Style:\\apolog.} & \makecell{Style:\\lecture} & \makecell{Style:\\redirect} & \makecell{Style:\\deflect} \\
\midrule
\multicolumn{10}{@{}l}{\textit{$<$1.3B}} \\
1 & Qwen2.5-0.5B & \TBD & \TBD & \TBD & \TBD & \TBD & \TBD & \TBD & \TBD \\
2 & Qwen3-0.6B & \TBD & \TBD & \TBD & \TBD & \TBD & \TBD & \TBD & \TBD \\
3 & Gemma-3-1B & \TBD & \TBD & \TBD & \TBD & \TBD & \TBD & \TBD & \TBD \\
4 & Llama-3.2-1B & \TBD & \TBD & \TBD & \TBD & \TBD & \TBD & \TBD & \TBD \\
\midrule
\multicolumn{10}{@{}l}{\textit{1.5--2.7B}} \\
5 & Qwen2.5-1.5B & \TBD & \TBD & \TBD & \TBD & \TBD & \TBD & \TBD & \TBD \\
6 & Qwen3-1.7B & \TBD & \TBD & \TBD & \TBD & \TBD & \TBD & \TBD & \TBD \\
7 & SmolLM2-1.7B & \TBD & \TBD & \TBD & \TBD & \TBD & \TBD & \TBD & \TBD \\
8 & Granite-3.3-2B & \TBD & \TBD & \TBD & \TBD & \TBD & \TBD & \TBD & \TBD \\
9 & Gemma-2-2B & \TBD & \TBD & \TBD & \TBD & \TBD & \TBD & \TBD & \TBD \\
10 & EXAONE-3.5-2.4B & \TBD & \TBD & \TBD & \TBD & \TBD & \TBD & \TBD & \TBD \\
\midrule
\multicolumn{10}{@{}l}{\textit{3--4.5B}} \\
11 & Qwen2.5-3B & \TBD & \TBD & \TBD & \TBD & \TBD & \TBD & \TBD & \TBD \\
12 & SmolLM3-3B & \TBD & \TBD & \TBD & \TBD & \TBD & \TBD & \TBD & \TBD \\
13 & Llama-3.2-3B & \TBD & \TBD & \TBD & \TBD & \TBD & \TBD & \TBD & \TBD \\
14 & Phi-3.5-mini & \TBD & \TBD & \TBD & \TBD & \TBD & \TBD & \TBD & \TBD \\
15 & Phi-4-mini & \TBD & \TBD & \TBD & \TBD & \TBD & \TBD & \TBD & \TBD \\
16 & Qwen3-4B-2507 & \TBD & \TBD & \TBD & \TBD & \TBD & \TBD & \TBD & \TBD \\
17 & Gemma-3-4B & \TBD & \TBD & \TBD & \TBD & \TBD & \TBD & \TBD & \TBD \\
18 & Gemma-3n-E4B & \TBD & \TBD & \TBD & \TBD & \TBD & \TBD & \TBD & \TBD \\
\midrule
\multicolumn{10}{@{}l}{\textit{7--9B}} \\
19 & Mistral-7B-v0.3 & \TBD & \TBD & \TBD & \TBD & \TBD & \TBD & \TBD & \TBD \\
20 & OLMo-2-7B & \TBD & \TBD & \TBD & \TBD & \TBD & \TBD & \TBD & \TBD \\
21 & Falcon3-7B & \TBD & \TBD & \TBD & \TBD & \TBD & \TBD & \TBD & \TBD \\
22 & Qwen2.5-7B & \TBD & \TBD & \TBD & \TBD & \TBD & \TBD & \TBD & \TBD \\
23 & Ministral-8B & \TBD & \TBD & \TBD & \TBD & \TBD & \TBD & \TBD & \TBD \\
24 & Llama-3.1-8B & \TBD & \TBD & \TBD & \TBD & \TBD & \TBD & \TBD & \TBD \\
25 & Hermes-3-8B & \TBD & \TBD & \TBD & \TBD & \TBD & \TBD & \TBD & \TBD \\
26 & R1-Distill-Llama-8B & \TBD & \TBD & \TBD & \TBD & \TBD & \TBD & \TBD & \TBD \\
27 & Aya-Expanse-8B & \TBD & \TBD & \TBD & \TBD & \TBD & \TBD & \TBD & \TBD \\
28 & Granite-3.3-8B & \TBD & \TBD & \TBD & \TBD & \TBD & \TBD & \TBD & \TBD \\
29 & Qwen3-8B & \TBD & \TBD & \TBD & \TBD & \TBD & \TBD & \TBD & \TBD \\
30 & Gemma-2-9B & \TBD & \TBD & \TBD & \TBD & \TBD & \TBD & \TBD & \TBD \\
\midrule
 & \textit{Mean (30 models)} & \TBD & \TBD & \TBD & \TBD & \TBD & \TBD & \TBD & \TBD \\
 & \textit{Spearman $\rho$ with $\log$ params} & \TBD & \TBD & \TBD & \TBD & \TBD & \TBD & \TBD & \TBD \\
\bottomrule
\end{tabular}
\end{table}


\begin{figure}[htbp]
\centering
\figplaceholder[0.25\textheight]{Scatter: $x$ = over-refusal rate on safe prompts, $y$ = compliance rate on unsafe prompts, one point per model, coloured by family. The ideal corner is bottom-left.}
\caption{The refusal--compliance trade-off across the 30 models.}
\label{fig:refusal-tradeoff}
\end{figure}

\begin{observation}[\Bx{4} refusal]
\questions{
\item Where do the models lie on the over-refusal vs.\ unsafe-compliance trade-off? Are there models that do well on both?
\item Which refusal style dominates per family? Do some families lecture more?
\item How do the community fine-tune (Hermes~3) and the reasoning distill compare with Llama-3.1-8B-Instruct?
}
\TBD
\end{observation}

% ---------------------------------------------------------------------
\subsection{Cluster C: Epistemic conduct}

\subsubsection{\Bx{5} Hedge faithfulness}
% Auto-generated skeleton: replace each \TBD with your measured value.
\begin{table}[htbp]
\centering
\caption{\Bx{5} Hedge faithfulness on TriviaQA. Accuracy of the greedy answer, mean intrinsic confidence (sampling consistency), mean decisiveness, cMFG, and ECE of lexicon-mapped verbal confidence.}
\label{tab:b5}
\scriptsize
\setlength{\tabcolsep}{3pt}
\begin{tabular}{@{}r l cccccc@{}}
\toprule
\# & Model & \makecell{Acc.} & \makecell{Mean\\confid.} & \makecell{Mean\\decisive.} & \makecell{cMFG\\$\uparrow$} & \makecell{ECE\\$\downarrow$} & \makecell{\% answers\\with hedge} \\
\midrule
\multicolumn{8}{@{}l}{\textit{$<$1.3B}} \\
1 & Qwen2.5-0.5B & \TBD & \TBD & \TBD & \TBD & \TBD & \TBD \\
2 & Qwen3-0.6B & \TBD & \TBD & \TBD & \TBD & \TBD & \TBD \\
3 & Gemma-3-1B & \TBD & \TBD & \TBD & \TBD & \TBD & \TBD \\
4 & Llama-3.2-1B & \TBD & \TBD & \TBD & \TBD & \TBD & \TBD \\
\midrule
\multicolumn{8}{@{}l}{\textit{1.5--2.7B}} \\
5 & Qwen2.5-1.5B & \TBD & \TBD & \TBD & \TBD & \TBD & \TBD \\
6 & Qwen3-1.7B & \TBD & \TBD & \TBD & \TBD & \TBD & \TBD \\
7 & SmolLM2-1.7B & \TBD & \TBD & \TBD & \TBD & \TBD & \TBD \\
8 & Granite-3.3-2B & \TBD & \TBD & \TBD & \TBD & \TBD & \TBD \\
9 & Gemma-2-2B & \TBD & \TBD & \TBD & \TBD & \TBD & \TBD \\
10 & EXAONE-3.5-2.4B & \TBD & \TBD & \TBD & \TBD & \TBD & \TBD \\
\midrule
\multicolumn{8}{@{}l}{\textit{3--4.5B}} \\
11 & Qwen2.5-3B & \TBD & \TBD & \TBD & \TBD & \TBD & \TBD \\
12 & SmolLM3-3B & \TBD & \TBD & \TBD & \TBD & \TBD & \TBD \\
13 & Llama-3.2-3B & \TBD & \TBD & \TBD & \TBD & \TBD & \TBD \\
14 & Phi-3.5-mini & \TBD & \TBD & \TBD & \TBD & \TBD & \TBD \\
15 & Phi-4-mini & \TBD & \TBD & \TBD & \TBD & \TBD & \TBD \\
16 & Qwen3-4B-2507 & \TBD & \TBD & \TBD & \TBD & \TBD & \TBD \\
17 & Gemma-3-4B & \TBD & \TBD & \TBD & \TBD & \TBD & \TBD \\
18 & Gemma-3n-E4B & \TBD & \TBD & \TBD & \TBD & \TBD & \TBD \\
\midrule
\multicolumn{8}{@{}l}{\textit{7--9B}} \\
19 & Mistral-7B-v0.3 & \TBD & \TBD & \TBD & \TBD & \TBD & \TBD \\
20 & OLMo-2-7B & \TBD & \TBD & \TBD & \TBD & \TBD & \TBD \\
21 & Falcon3-7B & \TBD & \TBD & \TBD & \TBD & \TBD & \TBD \\
22 & Qwen2.5-7B & \TBD & \TBD & \TBD & \TBD & \TBD & \TBD \\
23 & Ministral-8B & \TBD & \TBD & \TBD & \TBD & \TBD & \TBD \\
24 & Llama-3.1-8B & \TBD & \TBD & \TBD & \TBD & \TBD & \TBD \\
25 & Hermes-3-8B & \TBD & \TBD & \TBD & \TBD & \TBD & \TBD \\
26 & R1-Distill-Llama-8B & \TBD & \TBD & \TBD & \TBD & \TBD & \TBD \\
27 & Aya-Expanse-8B & \TBD & \TBD & \TBD & \TBD & \TBD & \TBD \\
28 & Granite-3.3-8B & \TBD & \TBD & \TBD & \TBD & \TBD & \TBD \\
29 & Qwen3-8B & \TBD & \TBD & \TBD & \TBD & \TBD & \TBD \\
30 & Gemma-2-9B & \TBD & \TBD & \TBD & \TBD & \TBD & \TBD \\
\midrule
 & \textit{Mean (30 models)} & \TBD & \TBD & \TBD & \TBD & \TBD & \TBD \\
 & \textit{Spearman $\rho$ with $\log$ params} & \TBD & \TBD & \TBD & \TBD & \TBD & \TBD \\
\bottomrule
\end{tabular}
\end{table}


\begin{figure}[htbp]
\centering
\figplaceholder[0.22\textheight]{Reliability-style plot: $x$ = intrinsic confidence (binned), $y$ = mean decisiveness; one line per size bin; the diagonal means faithful.}
\caption{Decisiveness vs.\ intrinsic confidence.}
\label{fig:hedging}
\end{figure}

\begin{observation}[\Bx{5} hedge faithfulness]
\questions{
\item Are small models over-decisive (decisive even when inconsistent)? Does faithfulness improve with size?
\item Do models hedge at all when asked? What is the hedge rate?
\item Does cMFG correlate with ECE? Does it correlate with sycophancy (\Bx{1})?
}
\TBD
\end{observation}

\subsubsection{\Bx{6} Clarification-seeking}
% Auto-generated skeleton: replace each \TBD with your measured value.
\begin{table}[htbp]
\centering
\caption{\Bx{6} Clarification-seeking on CLAMBER. Discrimination = clarification rate on ambiguous queries minus clarification rate on clear queries.}
\label{tab:b6}
\scriptsize
\setlength{\tabcolsep}{3pt}
\begin{tabular}{@{}r l ccccc@{}}
\toprule
\# & Model & \makecell{Clarify\\(ambig.)} & \makecell{Multi-\\answer} & \makecell{Direct\\answer} & \makecell{Clarify\\(clear)} & \makecell{Discrim.\\$\uparrow$} \\
\midrule
\multicolumn{7}{@{}l}{\textit{$<$1.3B}} \\
1 & Qwen2.5-0.5B & \TBD & \TBD & \TBD & \TBD & \TBD \\
2 & Qwen3-0.6B & \TBD & \TBD & \TBD & \TBD & \TBD \\
3 & Gemma-3-1B & \TBD & \TBD & \TBD & \TBD & \TBD \\
4 & Llama-3.2-1B & \TBD & \TBD & \TBD & \TBD & \TBD \\
\midrule
\multicolumn{7}{@{}l}{\textit{1.5--2.7B}} \\
5 & Qwen2.5-1.5B & \TBD & \TBD & \TBD & \TBD & \TBD \\
6 & Qwen3-1.7B & \TBD & \TBD & \TBD & \TBD & \TBD \\
7 & SmolLM2-1.7B & \TBD & \TBD & \TBD & \TBD & \TBD \\
8 & Granite-3.3-2B & \TBD & \TBD & \TBD & \TBD & \TBD \\
9 & Gemma-2-2B & \TBD & \TBD & \TBD & \TBD & \TBD \\
10 & EXAONE-3.5-2.4B & \TBD & \TBD & \TBD & \TBD & \TBD \\
\midrule
\multicolumn{7}{@{}l}{\textit{3--4.5B}} \\
11 & Qwen2.5-3B & \TBD & \TBD & \TBD & \TBD & \TBD \\
12 & SmolLM3-3B & \TBD & \TBD & \TBD & \TBD & \TBD \\
13 & Llama-3.2-3B & \TBD & \TBD & \TBD & \TBD & \TBD \\
14 & Phi-3.5-mini & \TBD & \TBD & \TBD & \TBD & \TBD \\
15 & Phi-4-mini & \TBD & \TBD & \TBD & \TBD & \TBD \\
16 & Qwen3-4B-2507 & \TBD & \TBD & \TBD & \TBD & \TBD \\
17 & Gemma-3-4B & \TBD & \TBD & \TBD & \TBD & \TBD \\
18 & Gemma-3n-E4B & \TBD & \TBD & \TBD & \TBD & \TBD \\
\midrule
\multicolumn{7}{@{}l}{\textit{7--9B}} \\
19 & Mistral-7B-v0.3 & \TBD & \TBD & \TBD & \TBD & \TBD \\
20 & OLMo-2-7B & \TBD & \TBD & \TBD & \TBD & \TBD \\
21 & Falcon3-7B & \TBD & \TBD & \TBD & \TBD & \TBD \\
22 & Qwen2.5-7B & \TBD & \TBD & \TBD & \TBD & \TBD \\
23 & Ministral-8B & \TBD & \TBD & \TBD & \TBD & \TBD \\
24 & Llama-3.1-8B & \TBD & \TBD & \TBD & \TBD & \TBD \\
25 & Hermes-3-8B & \TBD & \TBD & \TBD & \TBD & \TBD \\
26 & R1-Distill-Llama-8B & \TBD & \TBD & \TBD & \TBD & \TBD \\
27 & Aya-Expanse-8B & \TBD & \TBD & \TBD & \TBD & \TBD \\
28 & Granite-3.3-8B & \TBD & \TBD & \TBD & \TBD & \TBD \\
29 & Qwen3-8B & \TBD & \TBD & \TBD & \TBD & \TBD \\
30 & Gemma-2-9B & \TBD & \TBD & \TBD & \TBD & \TBD \\
\midrule
 & \textit{Mean (30 models)} & \TBD & \TBD & \TBD & \TBD & \TBD \\
 & \textit{Spearman $\rho$ with $\log$ params} & \TBD & \TBD & \TBD & \TBD & \TBD \\
\bottomrule
\end{tabular}
\end{table}


\begin{observation}[\Bx{6} clarification]
\questions{
\item How often do small models ask, compared with answering directly or giving multiple answers?
\item Is clarification discriminative (asked more for ambiguous than for clear queries), or indiscriminate?
\item Which CLAMBER ambiguity types trigger clarification, and which never do?
}
\TBD
\end{observation}

\subsubsection{\Bx{7} Cognitive biases}
% Auto-generated skeleton: replace each \TBD with your measured value.
\begin{table}[htbp]
\centering
\caption{\Bx{7} Cognitive-bias indices (0 = unbiased) and answer consistency across option permutations.}
\label{tab:b7}
\scriptsize
\setlength{\tabcolsep}{3pt}
\begin{tabular}{@{}r l ccccc@{}}
\toprule
\# & Model & \makecell{Anchoring\\index} & \makecell{Framing\\index} & \makecell{Position\\index} & \makecell{Most\\chosen pos.} & \makecell{Perm.\\consist.} \\
\midrule
\multicolumn{7}{@{}l}{\textit{$<$1.3B}} \\
1 & Qwen2.5-0.5B & \TBD & \TBD & \TBD & \TBD & \TBD \\
2 & Qwen3-0.6B & \TBD & \TBD & \TBD & \TBD & \TBD \\
3 & Gemma-3-1B & \TBD & \TBD & \TBD & \TBD & \TBD \\
4 & Llama-3.2-1B & \TBD & \TBD & \TBD & \TBD & \TBD \\
\midrule
\multicolumn{7}{@{}l}{\textit{1.5--2.7B}} \\
5 & Qwen2.5-1.5B & \TBD & \TBD & \TBD & \TBD & \TBD \\
6 & Qwen3-1.7B & \TBD & \TBD & \TBD & \TBD & \TBD \\
7 & SmolLM2-1.7B & \TBD & \TBD & \TBD & \TBD & \TBD \\
8 & Granite-3.3-2B & \TBD & \TBD & \TBD & \TBD & \TBD \\
9 & Gemma-2-2B & \TBD & \TBD & \TBD & \TBD & \TBD \\
10 & EXAONE-3.5-2.4B & \TBD & \TBD & \TBD & \TBD & \TBD \\
\midrule
\multicolumn{7}{@{}l}{\textit{3--4.5B}} \\
11 & Qwen2.5-3B & \TBD & \TBD & \TBD & \TBD & \TBD \\
12 & SmolLM3-3B & \TBD & \TBD & \TBD & \TBD & \TBD \\
13 & Llama-3.2-3B & \TBD & \TBD & \TBD & \TBD & \TBD \\
14 & Phi-3.5-mini & \TBD & \TBD & \TBD & \TBD & \TBD \\
15 & Phi-4-mini & \TBD & \TBD & \TBD & \TBD & \TBD \\
16 & Qwen3-4B-2507 & \TBD & \TBD & \TBD & \TBD & \TBD \\
17 & Gemma-3-4B & \TBD & \TBD & \TBD & \TBD & \TBD \\
18 & Gemma-3n-E4B & \TBD & \TBD & \TBD & \TBD & \TBD \\
\midrule
\multicolumn{7}{@{}l}{\textit{7--9B}} \\
19 & Mistral-7B-v0.3 & \TBD & \TBD & \TBD & \TBD & \TBD \\
20 & OLMo-2-7B & \TBD & \TBD & \TBD & \TBD & \TBD \\
21 & Falcon3-7B & \TBD & \TBD & \TBD & \TBD & \TBD \\
22 & Qwen2.5-7B & \TBD & \TBD & \TBD & \TBD & \TBD \\
23 & Ministral-8B & \TBD & \TBD & \TBD & \TBD & \TBD \\
24 & Llama-3.1-8B & \TBD & \TBD & \TBD & \TBD & \TBD \\
25 & Hermes-3-8B & \TBD & \TBD & \TBD & \TBD & \TBD \\
26 & R1-Distill-Llama-8B & \TBD & \TBD & \TBD & \TBD & \TBD \\
27 & Aya-Expanse-8B & \TBD & \TBD & \TBD & \TBD & \TBD \\
28 & Granite-3.3-8B & \TBD & \TBD & \TBD & \TBD & \TBD \\
29 & Qwen3-8B & \TBD & \TBD & \TBD & \TBD & \TBD \\
30 & Gemma-2-9B & \TBD & \TBD & \TBD & \TBD & \TBD \\
\midrule
 & \textit{Mean (30 models)} & \TBD & \TBD & \TBD & \TBD & \TBD \\
 & \textit{Spearman $\rho$ with $\log$ params} & \TBD & \TBD & \TBD & \TBD & \TBD \\
\bottomrule
\end{tabular}
\end{table}


\begin{observation}[\Bx{7} cognitive biases]
\questions{
\item Which models are anchored or framed most? Are the effects human-like in direction?
\item Is there a position preference (\eg for option A), and does it shrink with size?
\item Is low cognitive bias in some tiny models real, or an artefact of random or invalid answers?
}
\TBD
\end{observation}

% ---------------------------------------------------------------------
\subsection{Cluster D: Explicit vs.\ implicit social bias}
% Auto-generated skeleton: replace each \TBD with your measured value.
\begin{table}[htbp]
\centering
\caption{\Bx{8} Explicit vs.\ implicit social bias. BBQ bias score on ambiguous items, ``unknown'' rate, accuracy on disambiguated items, IAT-style association bias, and the implicit--explicit gap.}
\label{tab:b8}
\scriptsize
\setlength{\tabcolsep}{3pt}
\begin{tabular}{@{}r l ccccc@{}}
\toprule
\# & Model & \makecell{BBQ bias\\(ambig.)} & \makecell{``Unknown''\\rate} & \makecell{BBQ acc.\\(disamb.)} & \makecell{IAT\\bias} & \makecell{Gap\\(IAT$-$BBQ)} \\
\midrule
\multicolumn{7}{@{}l}{\textit{$<$1.3B}} \\
1 & Qwen2.5-0.5B & \TBD & \TBD & \TBD & \TBD & \TBD \\
2 & Qwen3-0.6B & \TBD & \TBD & \TBD & \TBD & \TBD \\
3 & Gemma-3-1B & \TBD & \TBD & \TBD & \TBD & \TBD \\
4 & Llama-3.2-1B & \TBD & \TBD & \TBD & \TBD & \TBD \\
\midrule
\multicolumn{7}{@{}l}{\textit{1.5--2.7B}} \\
5 & Qwen2.5-1.5B & \TBD & \TBD & \TBD & \TBD & \TBD \\
6 & Qwen3-1.7B & \TBD & \TBD & \TBD & \TBD & \TBD \\
7 & SmolLM2-1.7B & \TBD & \TBD & \TBD & \TBD & \TBD \\
8 & Granite-3.3-2B & \TBD & \TBD & \TBD & \TBD & \TBD \\
9 & Gemma-2-2B & \TBD & \TBD & \TBD & \TBD & \TBD \\
10 & EXAONE-3.5-2.4B & \TBD & \TBD & \TBD & \TBD & \TBD \\
\midrule
\multicolumn{7}{@{}l}{\textit{3--4.5B}} \\
11 & Qwen2.5-3B & \TBD & \TBD & \TBD & \TBD & \TBD \\
12 & SmolLM3-3B & \TBD & \TBD & \TBD & \TBD & \TBD \\
13 & Llama-3.2-3B & \TBD & \TBD & \TBD & \TBD & \TBD \\
14 & Phi-3.5-mini & \TBD & \TBD & \TBD & \TBD & \TBD \\
15 & Phi-4-mini & \TBD & \TBD & \TBD & \TBD & \TBD \\
16 & Qwen3-4B-2507 & \TBD & \TBD & \TBD & \TBD & \TBD \\
17 & Gemma-3-4B & \TBD & \TBD & \TBD & \TBD & \TBD \\
18 & Gemma-3n-E4B & \TBD & \TBD & \TBD & \TBD & \TBD \\
\midrule
\multicolumn{7}{@{}l}{\textit{7--9B}} \\
19 & Mistral-7B-v0.3 & \TBD & \TBD & \TBD & \TBD & \TBD \\
20 & OLMo-2-7B & \TBD & \TBD & \TBD & \TBD & \TBD \\
21 & Falcon3-7B & \TBD & \TBD & \TBD & \TBD & \TBD \\
22 & Qwen2.5-7B & \TBD & \TBD & \TBD & \TBD & \TBD \\
23 & Ministral-8B & \TBD & \TBD & \TBD & \TBD & \TBD \\
24 & Llama-3.1-8B & \TBD & \TBD & \TBD & \TBD & \TBD \\
25 & Hermes-3-8B & \TBD & \TBD & \TBD & \TBD & \TBD \\
26 & R1-Distill-Llama-8B & \TBD & \TBD & \TBD & \TBD & \TBD \\
27 & Aya-Expanse-8B & \TBD & \TBD & \TBD & \TBD & \TBD \\
28 & Granite-3.3-8B & \TBD & \TBD & \TBD & \TBD & \TBD \\
29 & Qwen3-8B & \TBD & \TBD & \TBD & \TBD & \TBD \\
30 & Gemma-2-9B & \TBD & \TBD & \TBD & \TBD & \TBD \\
\midrule
 & \textit{Mean (30 models)} & \TBD & \TBD & \TBD & \TBD & \TBD \\
 & \textit{Spearman $\rho$ with $\log$ params} & \TBD & \TBD & \TBD & \TBD & \TBD \\
\bottomrule
\end{tabular}
\end{table}


\begin{figure}[htbp]
\centering
\figplaceholder[0.25\textheight]{Scatter: $x$ = explicit bias (BBQ ambiguous bias score), $y$ = implicit association bias; one point per model; the identity line is shown.}
\caption{Explicit vs.\ implicit bias per model.}
\label{fig:bias}
\end{figure}

\begin{observation}[\Bx{8} explicit vs.\ implicit bias]
\questions{
\item Is \Hyp{5} supported, \ie do models with low BBQ bias still show strong implicit associations?
\item Which BBQ categories and IAT domains show the largest gaps?
\item Do the ``unknown'' rates differ systematically by family or size?
}
\TBD
\end{observation}

% ---------------------------------------------------------------------
\subsection{\Bx{9} Verbosity and style}
% Auto-generated skeleton: replace each \TBD with your measured value.
\begin{table}[htbp]
\centering
\caption{\Bx{9} Verbosity and style, averaged over all free-text generations.}
\label{tab:b9}
\scriptsize
\setlength{\tabcolsep}{3pt}
\begin{tabular}{@{}r l ccccc@{}}
\toprule
\# & Model & \makecell{Mean\\tokens} & \makecell{Median\\tokens} & \makecell{Hedges /\\100 tok} & \makecell{Apologies /\\100 tok} & \makecell{Markdown /\\100 tok} \\
\midrule
\multicolumn{7}{@{}l}{\textit{$<$1.3B}} \\
1 & Qwen2.5-0.5B & \TBD & \TBD & \TBD & \TBD & \TBD \\
2 & Qwen3-0.6B & \TBD & \TBD & \TBD & \TBD & \TBD \\
3 & Gemma-3-1B & \TBD & \TBD & \TBD & \TBD & \TBD \\
4 & Llama-3.2-1B & \TBD & \TBD & \TBD & \TBD & \TBD \\
\midrule
\multicolumn{7}{@{}l}{\textit{1.5--2.7B}} \\
5 & Qwen2.5-1.5B & \TBD & \TBD & \TBD & \TBD & \TBD \\
6 & Qwen3-1.7B & \TBD & \TBD & \TBD & \TBD & \TBD \\
7 & SmolLM2-1.7B & \TBD & \TBD & \TBD & \TBD & \TBD \\
8 & Granite-3.3-2B & \TBD & \TBD & \TBD & \TBD & \TBD \\
9 & Gemma-2-2B & \TBD & \TBD & \TBD & \TBD & \TBD \\
10 & EXAONE-3.5-2.4B & \TBD & \TBD & \TBD & \TBD & \TBD \\
\midrule
\multicolumn{7}{@{}l}{\textit{3--4.5B}} \\
11 & Qwen2.5-3B & \TBD & \TBD & \TBD & \TBD & \TBD \\
12 & SmolLM3-3B & \TBD & \TBD & \TBD & \TBD & \TBD \\
13 & Llama-3.2-3B & \TBD & \TBD & \TBD & \TBD & \TBD \\
14 & Phi-3.5-mini & \TBD & \TBD & \TBD & \TBD & \TBD \\
15 & Phi-4-mini & \TBD & \TBD & \TBD & \TBD & \TBD \\
16 & Qwen3-4B-2507 & \TBD & \TBD & \TBD & \TBD & \TBD \\
17 & Gemma-3-4B & \TBD & \TBD & \TBD & \TBD & \TBD \\
18 & Gemma-3n-E4B & \TBD & \TBD & \TBD & \TBD & \TBD \\
\midrule
\multicolumn{7}{@{}l}{\textit{7--9B}} \\
19 & Mistral-7B-v0.3 & \TBD & \TBD & \TBD & \TBD & \TBD \\
20 & OLMo-2-7B & \TBD & \TBD & \TBD & \TBD & \TBD \\
21 & Falcon3-7B & \TBD & \TBD & \TBD & \TBD & \TBD \\
22 & Qwen2.5-7B & \TBD & \TBD & \TBD & \TBD & \TBD \\
23 & Ministral-8B & \TBD & \TBD & \TBD & \TBD & \TBD \\
24 & Llama-3.1-8B & \TBD & \TBD & \TBD & \TBD & \TBD \\
25 & Hermes-3-8B & \TBD & \TBD & \TBD & \TBD & \TBD \\
26 & R1-Distill-Llama-8B & \TBD & \TBD & \TBD & \TBD & \TBD \\
27 & Aya-Expanse-8B & \TBD & \TBD & \TBD & \TBD & \TBD \\
28 & Granite-3.3-8B & \TBD & \TBD & \TBD & \TBD & \TBD \\
29 & Qwen3-8B & \TBD & \TBD & \TBD & \TBD & \TBD \\
30 & Gemma-2-9B & \TBD & \TBD & \TBD & \TBD & \TBD \\
\midrule
 & \textit{Mean (30 models)} & \TBD & \TBD & \TBD & \TBD & \TBD \\
 & \textit{Spearman $\rho$ with $\log$ params} & \TBD & \TBD & \TBD & \TBD & \TBD \\
\bottomrule
\end{tabular}
\end{table}


\begin{observation}[\Bx{9} verbosity]
\questions{
\item Which families are most verbose? Does verbosity correlate with judged behaviours (a judge-bias check)?
\item Do results for \Bx{2}--\Bx{4} change when verbosity is controlled for in \cref{eq:glmm}?
}
\TBD
\end{observation}

% ---------------------------------------------------------------------
\subsection{RQ2: Recipe vs.\ scale}
\label{sec:res-recipe}

\begin{table}[htbp]
\centering
\caption{Variance decomposition: likelihood-ratio tests comparing nested versions of
\cref{eq:glmm}. Each cell gives $\Delta$deviance (and $p$) when the factor is added.}
\label{tab:variance}
\small
\begin{tabular}{@{}lcccccccc@{}}
\toprule
Factor added & \Bx{1} & \Bx{2} & \Bx{3} & \Bx{4} & \Bx{5} & \Bx{6} & \Bx{7} & \Bx{8} \\
\midrule
$\log$ params           & \TBD & \TBD & \TBD & \TBD & \TBD & \TBD & \TBD & \TBD \\
+ family                & \TBD & \TBD & \TBD & \TBD & \TBD & \TBD & \TBD & \TBD \\
+ model type            & \TBD & \TBD & \TBD & \TBD & \TBD & \TBD & \TBD & \TBD \\
+ capability (MMLU)     & \TBD & \TBD & \TBD & \TBD & \TBD & \TBD & \TBD & \TBD \\
\bottomrule
\end{tabular}
\end{table}

\begin{table}[htbp]
\centering
\caption{Same-base triplet: three post-training recipes on Llama-3.1-8B. Values are
primary metrics; $^{*}$ marks a significant McNemar difference from Llama-3.1-8B-Instruct.}
\label{tab:triplet}
\footnotesize
\setlength{\tabcolsep}{4pt}
\begin{tabular}{@{}lccccccccc@{}}
\toprule
Model & MMLU & \Bx{1} & \Bx{2} & \Bx{3} & \Bx{4} & \Bx{5} & \Bx{6} & \Bx{7} & \Bx{8} \\
\midrule
Llama-3.1-8B-Instruct       & \TBD & \TBD & \TBD & \TBD & \TBD & \TBD & \TBD & \TBD & \TBD \\
Hermes-3-Llama-3.1-8B       & \TBD & \TBD & \TBD & \TBD & \TBD & \TBD & \TBD & \TBD & \TBD \\
DeepSeek-R1-Distill-Llama-8B& \TBD & \TBD & \TBD & \TBD & \TBD & \TBD & \TBD & \TBD & \TBD \\
\bottomrule
\end{tabular}
\end{table}

\begin{table}[htbp]
\centering
\caption{Generation pairs at matched size, and thinking on vs.\ off. Cells give
$\Delta$ = newer (or thinking-on) minus older (or thinking-off). $^{\dagger}$Nearest available size, not exactly matched.}
\label{tab:generation}
\footnotesize
\setlength{\tabcolsep}{4pt}
\begin{tabular}{@{}lccccccccc@{}}
\toprule
Contrast & MMLU & \Bx{1} & \Bx{2} & \Bx{3} & \Bx{4} & \Bx{5} & \Bx{6} & \Bx{7} & \Bx{8} \\
\midrule
Qwen2.5-0.5B $\rightarrow$ Qwen3-0.6B & \TBD & \TBD & \TBD & \TBD & \TBD & \TBD & \TBD & \TBD & \TBD \\
Qwen2.5-1.5B $\rightarrow$ Qwen3-1.7B & \TBD & \TBD & \TBD & \TBD & \TBD & \TBD & \TBD & \TBD & \TBD \\
Qwen2.5-7B $\rightarrow$ Qwen3-8B     & \TBD & \TBD & \TBD & \TBD & \TBD & \TBD & \TBD & \TBD & \TBD \\
Gemma-2-2B $\rightarrow$ Gemma-3-4B$^{\dagger}$& \TBD & \TBD & \TBD & \TBD & \TBD & \TBD & \TBD & \TBD & \TBD \\
Phi-3.5-mini $\rightarrow$ Phi-4-mini & \TBD & \TBD & \TBD & \TBD & \TBD & \TBD & \TBD & \TBD & \TBD \\
Mistral-7B $\rightarrow$ Ministral-8B & \TBD & \TBD & \TBD & \TBD & \TBD & \TBD & \TBD & \TBD & \TBD \\
\midrule
Qwen3-1.7B: think off $\rightarrow$ on & \TBD & \TBD & \TBD & \TBD & \TBD & \TBD & \TBD & \TBD & \TBD \\
SmolLM3-3B: think off $\rightarrow$ on & \TBD & \TBD & \TBD & \TBD & \TBD & \TBD & \TBD & \TBD & \TBD \\
Qwen3-8B: think off $\rightarrow$ on   & \TBD & \TBD & \TBD & \TBD & \TBD & \TBD & \TBD & \TBD & \TBD \\
\bottomrule
\end{tabular}
\end{table}

\begin{observation}[RQ2 --- recipe vs.\ scale]
\questions{
\item Does family or post-training explain more variance than size (\cref{tab:variance})? Is \Hyp{2} supported?
\item How different are the three Llama-3.1-8B variants? On which behaviours does the recipe matter most?
\item Did newer generations get less sycophantic and less over-refusing, or did they trade one for the other?
\item What does thinking mode change? Does reasoning reduce or increase sycophancy?
\item Do fully open models (OLMo~2, SmolLM) differ from closed-data models of similar size?
}
\TBD
\end{observation}

% ---------------------------------------------------------------------
\subsection{RQ3: Effects of quantization}
\label{sec:res-quant}

% Auto-generated skeleton: replace each \TBD with your measured value.
\begin{landscape}
\begin{table}[p]
\centering
\caption{Quantization sub-study: primary metric per behaviour at each precision level. F16 rows give absolute values; the other rows give $\Delta$ from F16 (percentage points), with $^{*}$ marking significance after BH correction.}
\label{tab:quant-aggregate}
\scriptsize
\setlength{\tabcolsep}{4pt}
\begin{tabular}{@{}llcccccccccc@{}}
\toprule
Model & Level & Eff.\ bpw & MMLU & B1 & B2 & B3 & B4 & B5 & B6 & B7 & B8 \\
\midrule
\multirow{5}{*}{Qwen3-1.7B} & F16 & \TBD & \TBD & \TBD & \TBD & \TBD & \TBD & \TBD & \TBD & \TBD & \TBD \\
 & Q8\_0 & \TBD & \TBD & \TBD & \TBD & \TBD & \TBD & \TBD & \TBD & \TBD & \TBD \\
 & Q4\_K\_M & \TBD & \TBD & \TBD & \TBD & \TBD & \TBD & \TBD & \TBD & \TBD & \TBD \\
 & Q3\_K\_M & \TBD & \TBD & \TBD & \TBD & \TBD & \TBD & \TBD & \TBD & \TBD & \TBD \\
 & Q2\_K & \TBD & \TBD & \TBD & \TBD & \TBD & \TBD & \TBD & \TBD & \TBD & \TBD \\
\midrule
\multirow{5}{*}{Qwen3-8B} & F16 & \TBD & \TBD & \TBD & \TBD & \TBD & \TBD & \TBD & \TBD & \TBD & \TBD \\
 & Q8\_0 & \TBD & \TBD & \TBD & \TBD & \TBD & \TBD & \TBD & \TBD & \TBD & \TBD \\
 & Q4\_K\_M & \TBD & \TBD & \TBD & \TBD & \TBD & \TBD & \TBD & \TBD & \TBD & \TBD \\
 & Q3\_K\_M & \TBD & \TBD & \TBD & \TBD & \TBD & \TBD & \TBD & \TBD & \TBD & \TBD \\
 & Q2\_K & \TBD & \TBD & \TBD & \TBD & \TBD & \TBD & \TBD & \TBD & \TBD & \TBD \\
\midrule
\multirow{5}{*}{Llama-3.2-3B} & F16 & \TBD & \TBD & \TBD & \TBD & \TBD & \TBD & \TBD & \TBD & \TBD & \TBD \\
 & Q8\_0 & \TBD & \TBD & \TBD & \TBD & \TBD & \TBD & \TBD & \TBD & \TBD & \TBD \\
 & Q4\_K\_M & \TBD & \TBD & \TBD & \TBD & \TBD & \TBD & \TBD & \TBD & \TBD & \TBD \\
 & Q3\_K\_M & \TBD & \TBD & \TBD & \TBD & \TBD & \TBD & \TBD & \TBD & \TBD & \TBD \\
 & Q2\_K & \TBD & \TBD & \TBD & \TBD & \TBD & \TBD & \TBD & \TBD & \TBD & \TBD \\
\midrule
\multirow{5}{*}{Llama-3.1-8B} & F16 & \TBD & \TBD & \TBD & \TBD & \TBD & \TBD & \TBD & \TBD & \TBD & \TBD \\
 & Q8\_0 & \TBD & \TBD & \TBD & \TBD & \TBD & \TBD & \TBD & \TBD & \TBD & \TBD \\
 & Q4\_K\_M & \TBD & \TBD & \TBD & \TBD & \TBD & \TBD & \TBD & \TBD & \TBD & \TBD \\
 & Q3\_K\_M & \TBD & \TBD & \TBD & \TBD & \TBD & \TBD & \TBD & \TBD & \TBD & \TBD \\
 & Q2\_K & \TBD & \TBD & \TBD & \TBD & \TBD & \TBD & \TBD & \TBD & \TBD & \TBD \\
\midrule
\multirow{5}{*}{Gemma-3-4B} & F16 & \TBD & \TBD & \TBD & \TBD & \TBD & \TBD & \TBD & \TBD & \TBD & \TBD \\
 & Q8\_0 & \TBD & \TBD & \TBD & \TBD & \TBD & \TBD & \TBD & \TBD & \TBD & \TBD \\
 & Q4\_K\_M & \TBD & \TBD & \TBD & \TBD & \TBD & \TBD & \TBD & \TBD & \TBD & \TBD \\
 & Q3\_K\_M & \TBD & \TBD & \TBD & \TBD & \TBD & \TBD & \TBD & \TBD & \TBD & \TBD \\
 & Q2\_K & \TBD & \TBD & \TBD & \TBD & \TBD & \TBD & \TBD & \TBD & \TBD & \TBD \\
\midrule
\multirow{5}{*}{Phi-4-mini} & F16 & \TBD & \TBD & \TBD & \TBD & \TBD & \TBD & \TBD & \TBD & \TBD & \TBD \\
 & Q8\_0 & \TBD & \TBD & \TBD & \TBD & \TBD & \TBD & \TBD & \TBD & \TBD & \TBD \\
 & Q4\_K\_M & \TBD & \TBD & \TBD & \TBD & \TBD & \TBD & \TBD & \TBD & \TBD & \TBD \\
 & Q3\_K\_M & \TBD & \TBD & \TBD & \TBD & \TBD & \TBD & \TBD & \TBD & \TBD & \TBD \\
 & Q2\_K & \TBD & \TBD & \TBD & \TBD & \TBD & \TBD & \TBD & \TBD & \TBD & \TBD \\
\bottomrule
\end{tabular}
\end{table}
\end{landscape}

% Auto-generated skeleton: replace each \TBD with your measured value.
\begin{table}[htbp]
\centering
\caption{Quantization sub-study: paired item-level flip rate (\%) relative to F16, averaged over the six models. In brackets: share of flips in the harmful direction. The final column gives MMLU flips for comparison.}
\label{tab:quant-flips}
\small
\setlength{\tabcolsep}{4pt}
\resizebox{\linewidth}{!}{\begin{tabular}{@{}lccccccccc@{}}
\toprule
Level & B1 & B2 & B3 & B4 & B5 & B6 & B7 & B8 & MMLU \\
\midrule
Q8\_0 & \TBD\,(\TBD) & \TBD\,(\TBD) & \TBD\,(\TBD) & \TBD\,(\TBD) & \TBD\,(\TBD) & \TBD\,(\TBD) & \TBD\,(\TBD) & \TBD\,(\TBD) & \TBD \\
Q4\_K\_M & \TBD\,(\TBD) & \TBD\,(\TBD) & \TBD\,(\TBD) & \TBD\,(\TBD) & \TBD\,(\TBD) & \TBD\,(\TBD) & \TBD\,(\TBD) & \TBD\,(\TBD) & \TBD \\
Q3\_K\_M & \TBD\,(\TBD) & \TBD\,(\TBD) & \TBD\,(\TBD) & \TBD\,(\TBD) & \TBD\,(\TBD) & \TBD\,(\TBD) & \TBD\,(\TBD) & \TBD\,(\TBD) & \TBD \\
Q2\_K & \TBD\,(\TBD) & \TBD\,(\TBD) & \TBD\,(\TBD) & \TBD\,(\TBD) & \TBD\,(\TBD) & \TBD\,(\TBD) & \TBD\,(\TBD) & \TBD\,(\TBD) & \TBD \\
\bottomrule
\end{tabular}}
\end{table}


\begin{figure}[htbp]
\centering
\figplaceholder{Dose--response: $x$ = effective bits per weight (log scale), $y$ = paired
flip rate vs.\ F16; one line per behaviour, with MMLU as a dashed reference; one panel
per model, or pooled with CIs.}
\caption{Behaviour drift as a function of effective bits per weight.}
\label{fig:quant-dose}
\end{figure}

\begin{observation}[RQ3 --- quantization]
\questions{
\item Is Q8\_0 indistinguishable from F16? At which level do aggregate scores first move, and at which level do item flips first exceed MMLU flips?
\item Do open-ended, judged behaviours (\Bx{2}, \Bx{3}, \Bx{4}) drift more than multiple-choice behaviours (\Bx{1}, \Bx{7}, \Bx{8}) at matched MMLU accuracy (\Hyp{3})?
\item Is drift directional (\eg more sycophantic, fewer refusals, more decisive) or symmetric noise?
\item Do flips concentrate on uncertain items, as \citet{hua2026uncertainty} found for bias?
\item Are small models more fragile to quantization than 8B models (bits $\times$ size interaction)?
}
\TBD
\end{observation}

% ---------------------------------------------------------------------
\subsection{RQ4: Structure across behaviours}
\label{sec:res-structure}

\begin{figure}[htbp]
\centering
\begin{subfigure}[t]{0.48\linewidth}
\centering
\figplaceholder[0.22\textheight]{Spearman correlation matrix of primary metrics and MMLU/IFEval across 30 models (clustered).}
\caption{Correlations.}
\label{fig:corr}
\end{subfigure}\hfill
\begin{subfigure}[t]{0.48\linewidth}
\centering
\figplaceholder[0.22\textheight]{PCA biplot: models as points (coloured by family), behaviours as loading arrows.}
\caption{PCA.}
\label{fig:pca}
\end{subfigure}
\caption{Cross-behaviour structure.}
\label{fig:structure}
\end{figure}

\begin{observation}[RQ4 --- structure]
\questions{
\item Do behaviours cluster as hypothesised in \Hyp{4} (agreeableness vs.\ caution)?
\item How strongly is each behaviour correlated with MMLU (\Hyp{6})? Does capability predict user-facing behaviour?
\item Are there clear trade-offs (\eg low over-refusal together with high dark patterns)?
}
\TBD
\end{observation}

% ---------------------------------------------------------------------
\subsection{Summary of hypotheses}
\begin{table}[htbp]
\centering
\caption{Outcome of the pre-registered hypotheses.}
\label{tab:hypotheses}
\small
\begin{tabularx}{\textwidth}{@{}l>{\raggedright\arraybackslash}X l@{}}
\toprule
Hypothesis & Evidence & Verdict \\
\midrule
\Hyp{1} Size effects differ by behaviour & \TBD & \TBD \\
\Hyp{2} Recipe $>$ size & \TBD & \TBD \\
\Hyp{3} Quantization drift is hidden and concentrated in open-ended behaviours & \TBD & \TBD \\
\Hyp{4} Agreeableness and caution clusters & \TBD & \TBD \\
\Hyp{5} Implicit bias persists despite low explicit bias & \TBD & \TBD \\
\Hyp{6} Weak link to capability & \TBD & \TBD \\
\bottomrule
\end{tabularx}
\end{table}
